%% ---------------------------------------------------------------------------
\section{Background}
\label{sec:background}

\subsection{Content credentials}

A C2PA manifest is a record embedded in a media file that states who produced
it, with what device or software, and what was done to it since~\cite{c2pa25}.
Its assertions include a hard binding, a cryptographic hash over the content
that ties the claim to specific bytes. A signature over the claim makes the
record verifiable, and any change to the content breaks the binding.

The signature is a COSE \texttt{Sign1} structure~\cite{rfc9052} with a
detached payload. Its protected header carries the algorithm identifier and the
signer's certificates in the \texttt{x5chain} header of
RFC~9360~\cite{rfc9360}. C2PA narrows RFC~9360 in two ways that matter here.
The claim generator ``shall include the signer's certificate and all intermediate certificate authorities'' in \texttt{x5chain}, and ``the trust anchor's certificate (also called the root certificate) should not be included''; from version~2.5 it ``shall use only the integer 33'' as the label, in the protected bucket, while validators still accept the older string label
\texttt{x5chain} in either bucket~\cite{c2pa25}. A validator trusts a signer
whose certificate path ends at an anchor on the C2PA trust list, a list of
certificates that may be roots or issuing CAs~\cite{c2patrustlist}; a separate
list does the same for time-stamping authorities. In a JPEG the manifest
travels in APP11 segments, so the cost of a credential is the growth in file
size.

Version~2.4 of the specification, on which the measurements were planned,
permitted seven classical algorithms. Version~2.5, published in October 2026,
added ML-DSA-44, 65 and 87 under the COSE identifiers $-48$, $-49$ and
$-50$~\cite{c2pa25}. It adds nothing for SLH-DSA or Falcon.

\subsection{Post-quantum signatures}

ML-DSA~\cite{fips204} is the module-lattice signature scheme NIST
standardised in 2024, at three parameter sets. ML-DSA-44 targets NIST security
category~2 rather than category~1, so it is not the drop-in counterpart of a
128-bit classical parameter set that its name suggests. SLH-DSA~\cite{fips205}
is a hash-based alternative with conservative assumptions and much larger
signatures. FN-DSA, the standard NIST plans for Falcon as FIPS~206, had not
appeared even in draft by 23~October 2026; in September 2026 NIST asked for
comment on a revised plan that specifies fixed-point arithmetic for key
generation and signing~\cite{nist_fips206_status,nist_fips206_forum,%
nist_fips206_plan}. Composite signatures pair ML-DSA with a classical scheme: a verifier accepts only if both component signatures verify, so a forger must break both, and the signature stays secure as long as either algorithm does; the IETF specification is close to
publication~\cite{ounsworth2026composite}, and no COSE identifiers have been
assigned for them.

Signature and public-key sizes are what matter here, because both enter a
manifest: the signature once, and the public key and the issuer's signature
once per certificate. ML-DSA-44, 65 and 87 produce signatures of 2,420, 3,309
and 4,627 bytes and public keys of 1,312, 1,952 and 2,592, against 64 and 32
for Ed25519 and 256 and about 270 for RSA-2048. Certificates for ML-DSA and
SLH-DSA are specified by RFCs~9881 and~9909~\cite{rfc9881,rfc9909}, and CMS
signatures, which a timestamp token is, by RFCs~9882 and~9814~\cite{rfc9882,%
rfc9814}. The IANA COSE registry has no entry for FN-DSA or SLH-DSA; the
working-group draft for FN-DSA requested $-54$ and $-55$, but it expired on
16~September 2026 with both values unassigned~\cite{iana_cose,coseFalcon}.

\subsection{Carrying certificates, or not}

RFC~9360 gives a COSE structure two ways to identify its
signer~\cite{rfc9360}: \texttt{x5chain} carries certificates and \texttt{x5t}
carries a thumbprint, leaving the verifier to find the certificate elsewhere.
Within \texttt{x5chain} there is a further choice of how many certificates to
carry. Section~\ref{sec:carrying} measures four forms: leaf, intermediate and
root; leaf and intermediate, the form C2PA prescribes; the leaf alone; and the
thumbprint.
